% Generated draft: requires mathematical and source review.
\documentclass[11pt]{article}
\usepackage[margin=25mm]{geometry}
\usepackage{fontspec}
\setmainfont{Latin Modern Roman}
\usepackage{amsmath,amssymb,mathtools,unicode-math}
\setmathfont{Latin Modern Math}
\usepackage{enumitem,needspace,xurl,hyperref}
\setlength{\parindent}{0pt}
\setlength{\parskip}{5pt}
\setlength{\emergencystretch}{3em}
\newcommand{\N}{\mathbb{N}}
\newcommand{\Z}{\mathbb{Z}}
\newcommand{\Q}{\mathbb{Q}}
\newcommand{\R}{\mathbb{R}}
\newcommand{\C}{\mathbb{C}}
\newcommand{\F}{\mathbb{F}}
\newcommand{\A}{\mathbb{A}}
\newcommand{\PP}{\mathbb{P}}
\newcommand{\E}{\mathbb{E}}
\newcommand{\eps}{\varepsilon}
\newcommand{\dd}{\,\mathrm{d}}
\begin{document}
\section*{UnsolvedMath ID: 10400086}
\textbf{Source problem code:} catalogue entry
\subsection{Definitions and mathematical statement}
For every oriented \(3\)-manifold \(M\), Problem 4.8 asks to calculate the source's module \(\mathcal S_3(M)\) and to find a convenient methodology for doing so. Its defining coefficient ring, parameters and skein relations are not included in the supplied record.

A common interpretation of \(\mathcal S_3\) is the oriented HOMFLY-type skein module: take the free module over a commutative ring \(R\) on oriented link isotopy classes and quotient by a prescribed three-term relation. In one convention, with parameters \(v\) and \(z\) and \(v\) invertible, it is
\[
 v^{-1}[L_+]-v[L_-]=z[L_0],
\]
where the links agree outside a ball and differ there by a positive crossing, negative crossing and oriented smoothing. Framing and trivial-component normalization may be part of the chosen version. This conditional description explains the notation but does not fix the source module: different coefficient rings, parameter specializations and framing conventions must not be conflated.

Once the source definition is recovered, the task is to determine \(\mathcal S_3(M)\) as a module for each oriented \(M\), by a usable explicit description or presentation, and to develop a method for producing such descriptions. The record does not limit \(M\) to closed or compact manifolds and does not prescribe an input model, output format, complexity bound, or additional algebra structure.
\subsection{Short English statement}
Compute \(\mathcal S_3\) for every oriented three-manifold and provide a convenient method of calculation. The notation usually refers to an oriented three-term skein module, but its exact relations and coefficients are missing here. The supplied open-triage summary reports no clear resolution.
\subsection{Sources}
\begin{itemize}
\item Local UnsolvedMath registry, permanent ID \(10400086\), exact source code \texttt{catalogue entry}, Problem 4.8. No individual proposer or proposed year is supplied.
\item Ohtsuki, \emph{Problems on invariants of knots and 3-manifolds} (2002), Problem 4.8, PDF page 76: \url{https://msp.org/gtm/2002/04/gtm-2002-04-024s.pdf}. Bold printed headings distinguish 231 numbered items from references in surrounding exposition.
\item Extraction pdf-font-structure, accessed 2026-07-29. Source evidence NEEDS\_REVIEW presents an open problem and requires a current primary-literature audit. Research classification OPEN-TRIAGE; the supplied summary says calculation for all \(3\)-manifolds is not obviously resolved. No further specific references are supplied.
\item Dataset attribution: Ulam AI / UnsolvedMath Contributors, CC BY 4.0. Source rights NEEDS\_REVIEW: publicly accessible author/publisher copy, redistribution terms requiring review.
\end{itemize}
\end{document}
